The first conceptual inversion consists in conceiving the vacuum as a positive relational response.\par

In HMT, the vacuum is the minimal dynamical response of the system when its
two sheets must coexist, propagate, and close their difference while
preserving the orientation that distinguishes them.\par

The two sheets constitute two conjugate readings of the same angular antecedent. One preserves an orientation, and the other preserves the opposite orientation. Hence two channels, \(q_+\) and \(q_-\), appear as the two correlated spectral faces of a single state.\par

The operator\par

\[
\mathcal V_{90,120}
=
(I-T^{90})(I-T^{120})^{-1}
\]

compares two propagation depths: \(90\) and \(120\). Those depths are determined by the same elementary mode:\par

\[
90=3\cdot30,
\qquad
120=4\cdot30.
\]

Therefore, the response of each sheet can be rewritten as\par

\[
\frac{1-q^{90}}{1-q^{120}}
=
\frac{1+s+s^2}{1+s+s^2+s^3},
\qquad
s=q^{30}.
\]

And here is where the interpretation becomes genuinely powerful. The numerator contains three states of the mode \(30\); the denominator contains those same three states and one more. The vacuum is, therefore, an exact completion \(3\to4\).\par

That fourth term integrates and preserves the preceding three. It completes them. What is added is recorded as a positive completion residue:\par

\[
1-\frac{1+s+s^2}{1+s+s^2+s^3}
=
\frac{s^3}{1+s+s^2+s^3}>0.
\]

Vacuum memory is precisely the exact increment that allows three terms to become four. It is the positive and necessary trace of the completion.\par

From the two eigenvalues of that single operator arise the four constitutive relations:\par

\[
\widehat\varepsilon=r_+^2,
\qquad
\widehat\mu=r_-^2,
\qquad
\widehat Z=\frac{r_-}{r_+},
\qquad
\widehat c=\frac1{r_+r_-}.
\]

This completely changes the interpretation of permeability and permittivity.\par

Permittivity is the quadratic response of one sheet. Permeability is the quadratic response of the conjugate sheet. Impedance preserves the oriented relation between them. Propagation speed depends on their common product.\par

Thus, product and quotient perform two different tasks:\par

\begin{itemize}[leftmargin=*,itemsep=0.35em]
\item the product forgets which sheet was which and preserves the common propagation;
\item the quotient preserves what the relative orientation between them was.
\end{itemize}

The product answers the question: “what do the two leaves share?”. The quotient answers the question: “how are they distinguished?”.\par

Por eso\par

\[
\widehat\mu\,\widehat\varepsilon=\widehat c^{-2}
\]

y\par

\[
\frac{\widehat\mu}{\widehat\varepsilon}
=
\widehat Z^2
\]

are the two complementary ways of reading the same spectral pair.\par

The holographic structure appears here in a particularly transparent form: four visible magnitudes actually contain only two internal degrees of freedom. Any of certain suitable pairs allows the two sheets to be reconstructed. The constitutive publication preserves genealogy when product and quotient remain jointly available.\par

This is the deep repair of the ablation of \(\mu\) and \(\varepsilon\). The repair reunites their values with their common operator-theoretic origin.\par

The International System chart may subsequently transform these intrinsic sections into \(\mu_0\), \(\varepsilon_0\), \(Z_0\), and \(c\). Its function is to express, within a metrological convention, a structure whose provenance has already been selected by the operator.\par

The vacuum thereby ceases to be a repository of four constantes. becomes a unique spectral operation With two conjugate memories.\par

And that view has an important philosophical consequence: the vacuum is a compatibility relation between two oriented ways of realizing one and the same state.\par

